\documentclass{article}
\usepackage{tcee-repaso}
\tcetema[Transmisión monetaria]{3.A.37}{Política monetaria (II). Los mecanismos de transmisión de la política monetaria convencional. Rigideces de precios, rigideces de salarios y fricciones financieras. Los mecanismos de transmisión de la política monetaria no convencional.}
\tcefecha{05/11/2026}
\begin{document}
\cabecerarepaso
\begin{columnas}

% ==========================================================================
\bloque{Esquema del cante}
\esquemacante{%
\apartadocante{Introducción: «las cosas de palacio van despacio»}{2}
\apartadocante{I. Transmisión de la política convencional}{9}
\subcante{1. ¿Por qué tiene efectos reales el dinero?}{2,5}
\subcante{2. Los canales de transmisión}{3}
\subcante{3. Zona del euro y España}{1,5}
\subcante{4. Evidencia: efectos, retardos, identificación}{2}
\apartadocante{II. Rigideces y fricciones financieras}{9}
\subcante{1. Rigideces de precios: modelo neokeynesiano}{3}
\subcante{2. Rigideces de salarios: contratos escalonados}{2}
\subcante{3. Fricciones financieras: canales del crédito}{2,5}
\subcante{4. Hogares distintos: canal redistributivo}{1,5}
\apartadocante{III. Transmisión de la política no convencional}{9}
\subcante{1. El límite inferior de los tipos}{1,5}
\subcante{2. El modelo neokeynesiano en el límite}{7,5}
\apartadocante{Conclusión}{1}}
\emph{Pizarra}: \emph{izda.} (\hyperref[rep:I]{I}), escuelas y esquema de canales, que se queda; joroba en un rincón. \emph{Centro} (\hyperref[rep:II]{II}): 3 ecuaciones y gráficos que se borran. \emph{Dcha.} (\hyperref[rep:III]{III}): el modelo en el límite y su esquema de salidas. \emph{Ejemplo}: familia con hipoteca variable y pyme de azulejos.

% ==========================================================================
\bloque[2]{Introducción}
\begin{itemize}
\item \textbf{Enganche}: «Las cosas de palacio van despacio». 21/7/2022: BCE sube tipos por 1.ª vez en 11 años. Euríbor 12 m: $-0{,}50$\,\% (dic.\ 2021) $\to$ 4,16\,\% (oct.\ 2023); inflación zona del euro: techo 10,6\,\% (oct.\ 2022), cerca del 2\,\% solo en otoño 2024. Efecto máximo en precios: 29 meses. Y se repite: inflación 4\,\% (oct.\ 2026), subidas en jun.\ y sept.: depósito 2,50\,\%. \emph{¿Por qué caminos y con qué retraso llega el tipo a un día al gasto y a los precios, y qué hacer cuando ya no puede bajar?}
\item \textbf{Relevancia}: \emph{política}: se decide hoy sobre la economía de dentro de 1-2 años; sin canales ni retardos no se sabe cuánto mover ni cuándo parar. \emph{España} (bancos, hipotecas variables): antes y más fuerte \vertema{3.B.43}. \emph{Teoría}: efectos reales exigen fricciones.
\item \textbf{Contextualización}: \concepto{mecanismos de transmisión}: caminos de una decisión del BC (tipo, balance, comunicación) a precios de activos y financiación $\to$ demanda $\to$ $Y$ y $\pi$; cada camino, un \concepto{canal}. \concepto{Convencional}: tipo a muy CP (BCE desde 2024: facilidad de depósito, guía el €STR); \concepto{no convencional}: balance (compras, préstamos largos) y comunicación en el límite o si la transmisión se rompe. \concepto{Neutral}: solo mueve lo nominal (casi a LP, no a CP). \concepto{Rigidez nominal}: precios/salarios nominales no se ajustan enseguida. \concepto{Rigidez real}: el precio/salario \emph{relativo} deseado responde poco a la actividad; sola no da efectos reales, amplía los de la nominal (\autor{Ball}-\autor{Romer}, 1990) \vertema{3.A.7}. \concepto{Fricción financiera}: información o incentivos entre prestamista y prestatario $\Rightarrow$ coste del crédito depende del patrimonio.
\item \textbf{Problemática}: ¿por qué efectos reales y solo a CP? ¿Por qué canales y con qué retraso? ¿Las fricciones crean o amplifican? ¿Qué cambia si el tipo no puede bajar?
\item \textbf{Estructura}: \hyperref[rep:I]{I}, canales de la convencional; \hyperref[rep:II]{II}, rigideces y fricciones que la hacen posible y la amplifican; \hyperref[rep:III]{III}, canales de la no convencional.
\end{itemize}

% ==========================================================================
\bloque[9]{I. Transmisión de la política convencional\label{rep:I}}
\begin{ideaclave}
Efectos reales a CP, casi ninguno a LP. Llega por varios canales a la vez (tipo de interés, activos, tipo de cambio, crédito, expectativas) y despacio: más de 1 año en $Y$, más de 2 en $\pi$. Medirlo exige separar lo que decide el BC de su respuesta a la economía.
\end{ideaclave}

\sub{I.1. ¿Por qué tiene efectos reales el dinero?\label{rep:I.1}}
\begin{itemize}
\item \textbf{Clásicos}: \concepto{identidad contable del gasto} (\autor{Hume}, 1752): $MV\equiv PT$, siempre se cumple y no dice qué causa qué. \concepto{Teoría cuantitativa} (\autor{Fisher}, 1911): la identidad + 3 supuestos ($V$ estable; $T\approx Y$ de pleno empleo, $M$ exógena; causalidad $M\to P$) $\Rightarrow P=M\overline V/\overline Y$. En tasas, la identidad: $\gamma_M+\gamma_V\equiv\pi+\gamma_Y$; la teoría ($\gamma_V=0$): $\pi=\gamma_M-\gamma_Y$. Transmisión \emph{directa} por saldos reales; \concepto{dicotomía clásica}.
\item \textbf{\autor{Keynes}} (1936): \emph{indirecta}, \concepto{canal del tipo de interés}: $M$↑ $\Rightarrow$ $i$↓ $\Rightarrow$ $I$↑ $\Rightarrow$ multiplicador. 2 eslabones débiles: \emph{trampa de la liquidez} (el tipo no baja) e \emph{inversión insensible} (\emph{animal spirits}). Postkeynesianos (\autor{Kaldor}, 1982; \autor{Moore}, 1988): dinero endógeno, oferta horizontal al tipo oficial; los préstamos crean depósitos (\autor{McLeay}-\autor{Radia}-\autor{Thomas}, 2014).
\item \textbf{Síntesis}: IS-LM (\autor{Hicks}, 1937) \vertema{3.A.5}: $Y=(\overline A-bi)/(1-c)$, $M/P=kY-hi$ $\Rightarrow$ $\frac{\mathrm dY}{\mathrm d(M/P)}=\frac{b}{h(1-c)+bk}$: más eficaz con $h$↓ (LM inclinada) y $b$↑ (IS plana); eslabones débiles: $h\to\infty$, $b\to0$. \autor{Phillips} como menú (\autor{Samuelson}-\autor{Solow}, 1960).
\item \textbf{Monetaristas}: \emph{money matters}; \emph{reajuste de carteras} (bonos, acciones, capital, duraderos) + riqueza (\autor{Brunner}-\autor{Meltzer}, 1972). Retardos largos y variables: 6-9 meses a $PY$ y otros tantos a $P$ (\autor{Friedman}, 1961, 1970) $\Rightarrow$ regla del $k$\,\%. LP: \autor{Phillips} vertical.
\item \textbf{NMC}: solo lo \emph{no anticipado}. Islas (\autor{Lucas}, 1972): $\omega=\sigma^2_z/(\sigma^2_z+\sigma^2_p)$, $y_t=y^n+\psi(p_t-p^e_t)$ \vertema{3.A.6}; \concepto{proposición de inefectividad} (\autor{Sargent}-\autor{Wallace}, 1975): $y_t=y^n+(\psi\varepsilon_t+u_t)/(1+\psi)$. \autor{Ball}-\autor{Mankiw}-\autor{Romer} (1988), 43 países: la pendiente la explica el \emph{nivel} de inflación. Ciclo real: neutral incluso a CP.
\item \textbf{NEK}: expectativas racionales sin flexibilidad: con contratos escalonados también la política anticipada tiene efectos (\autor{Fischer}, 1977; \autor{Taylor}, 1980). Lo que daba la inefectividad era la flexibilidad.
\end{itemize}

\sub{I.2. Los canales de transmisión\label{rep:I.2}}
\begin{itemize}
\item \emph{1.er eslabón}: tipo oficial $\to$ interbancario $\to$ curva y tipos bancarios: $i^{(n)}_t=\frac1n\sum_k\mathbb E_t[i_{t+k}]+$ prima por plazo \vertema{3.B.24}: pesa lo que el BC anuncia.
\end{itemize}
\graficorepaso[0.62\linewidth]{canales-transmision}{Canales: del tipo oficial a la demanda, $Y$ y $\pi$ (1-2 años de retraso)}
\begin{itemize}
\item \textbf{Tipo de interés}: $r\approx i-\mathbb E_t\pi_{t+1}$; con precios rígidos, $i$↑ = $r$↑: inversión, $F_K(K^*)=r+\delta$ (\autor{Jorgenson}) \vertema{3.A.34}; consumo, \autor{Euler} $\Delta\ln C_{t+1}\approx\frac1\theta(r_t-\rho)$ (EIS baja) \vertema{3.A.33}; vivienda y duraderos, lo más sensible. $i\uparrow\Rightarrow r\uparrow\Rightarrow I,C^{\mathrm{dur}},C\downarrow\Rightarrow Y\downarrow\Rightarrow\pi\downarrow$. \emph{Familia}: 150.000 € a 25 años, Euríbor + 1: cuota $\approx$530 € (dic.\ 2021) $\to$ $\approx$890 € (oct.\ 2023): 360 € menos al mes, por la cuota, no por preferir ahorrar.
\item \textbf{Precios de activos}: $P^A=D/(i+\text{prima}-\gamma_D)$↓; \concepto{$q$ de \autor{Tobin}} (1969) $=P^A/(P_KK)$; $q<1$ no compensa invertir; cuenta la $q$ marginal, la bolsa da la media: coinciden con rend.\ constantes y competencia perfecta (\autor{Hayashi}, 1982). \emph{Efecto riqueza}: $C=c_YY^L+c_WW$ (ciclo vital) \vertema{3.A.33}.
\item \textbf{Tipo de cambio}: paridad no cubierta $i=i^*+\gamma E^e$ ($E$ €/\$): BCE sube y la Fed no $\Rightarrow$ € se aprecia hoy hasta esperar depreciación \vertema{3.B.14}. Efectos: $e=EP^*/P$↓ $\Rightarrow$ $X-M$↓ (\autor{Marshall}-\autor{Lerner}) \vertema{3.B.12}; inflación importada↓, traspaso incompleto, depende de la moneda de facturación (dólar: \autor{Gopinath}, 2020); menor en zona del euro que en cada país. \emph{Pyme}: azulejos más caros en EE.\,UU., menos pedidos. Hacia fuera: \concepto{ciclo financiero global} (\autor{Rey}, 2013; \autor{Miranda-Agrippino}-\autor{Rey}, 2020): cuando EE.\,UU.\ estornuda, el mundo se resfría.
\item \textbf{Crédito} (\autor{Bernanke}-\autor{Gertler}, 1995): cambia \emph{cuánto} crédito y a quién. $i^{\text{préstamo}}=i+s(\mathrm{RN})$. \concepto{Crédito bancario}: bancos dependientes de depósitos prestan menos; pymes sin sustituto. \concepto{Balance}: $i$↑ $\Rightarrow$ flujo de caja y garantías↓ $\Rightarrow$ prima↑. \concepto{Toma de riesgos} (\autor{Borio}-\autor{Zhu}, 2012): tipos bajos largo tiempo $\Rightarrow$ búsqueda de rentabilidad; España, registro de crédito: bancos con menos capital prestan a empresas más arriesgadas (\autor{Jiménez} et al., 2014).
\item \textbf{Expectativas}: curva de tipos y NKPC ($\pi_t=\beta\mathbb E_t\pi_{t+1}+\kappa x_t$): subida creíble $\Rightarrow$ $\pi^e$↓ $\Rightarrow$ salarios y precios fijados con el objetivo: desinflación con menos pérdida de $Y$ \vertema{3.A.36}. Clave del NEK y del límite inferior.
\item \emph{Retardos} (revisión gradual, inversión planificada, deudas escalonadas) $\Rightarrow$ el BC mira a la $\pi$ prevista a 1-2 años.
\end{itemize}
\graficorepaso[\linewidth]{cadenas-canales}{La cadena de cada canal: de $i$↑ a $Y$↓ y $\pi$↓}

\sub{I.3. La transmisión en la zona del euro y en España\label{rep:I.3}}
\begin{itemize}
\item Empresas más dependientes de bancos que en EE.\,UU. (\autor{De Fiore}-\autor{Uhlig}, 2011) $\Rightarrow$ pesa el crédito. 
\item \emph{2022-2023}: depósito $-0{,}50\to4$\,\% (450 pb): crédito nuevo a empresas +350 pb; hipotecas nuevas algo más de 200; plazo hogares $\approx$300 (zona) y 260 (España); vista 0 $\to$ 0,4\,\% (zona), 0,2\,\% (España). 
\item \emph{Tipo variable}: hipotecas vivas España 1,2\,\% (jun.\ 2022) $\to$ 3,7\,\% (feb.\ 2024); zona 1,6 $\to$ 2,5\,\%: canal del balance más rápido aquí; las nuevas, ya a tipo fijo o mixto. 
\item \emph{Bajadas} jun.\ 2024-jun.\ 2025: $-200$ pb; empresas $-150$, hipotecas $<50$.
\item Una política única: si las primas de un país se disparan, la transmisión se rompe (OMT, TPI) \vertema{3.B.43}.
\end{itemize}

\sub{I.4. Evidencia: efectos, retardos, identificación\label{rep:I.4}}
\begin{itemize}
\item \emph{Identificación}: el BC reacciona a la economía (sube si prevé $\pi$): aislar la \concepto{perturbación monetaria} ($\varepsilon_t$ de la regla de \autor{Taylor}).
\item \emph{Métodos}: \emph{historia} (\autor{Friedman}-\autor{Schwartz}, 1963) y \emph{San Luis} (\autor{Andersen}-\autor{Jordan}, 1968: \$1 $\to$ +5,8 de PNB nominal; pero el dinero responde a la renta). \emph{Narrativa}: \autor{Romer}-\autor{Romer} (1989, 6 episodios, el último \autor{Volcker} 1979; 2004, limpian de previsiones). \emph{VAR} (\autor{Sims}, 1992; \autor{Christiano}-\autor{Eichenbaum}-\autor{Evans}, 1996): $Y$ y $P$ no reaccionan en el periodo. \emph{Alta frecuencia} (30 min alrededor del anuncio): \autor{Gertler}-\autor{Karadi} (2015, EE.\,UU.); zona del euro, \autor{Altavilla} et al.\ (2019), \autor{Jarociński}-\autor{Karadi} (2020: separan sorpresa de política y de \emph{información}).
\item \emph{Resultados}: LP $\gamma_M$ y $\pi$ casi 1 a 1, $Y$ independiente. CP: $Y$ cae en \concepto{joroba}, fondo tras varios trimestres; narrativa: producción industrial $-4{,}3$\,\% $\approx$2 años tras +1 punto. Precios: bajan tras 22 meses (\autor{Romer}-\autor{Romer}); máximo a 29 meses de media (\autor{Havranek}-\autor{Rusnak}).
\item \concepto{Rompecabezas de los precios} (\emph{price puzzle}): en los 1.os VAR los precios \emph{subían} (\autor{Sims}, 1992): el BC sube porque prevé $\pi$. Se corrige con sus previsiones y separando las sorpresas de información.
\end{itemize}
\graficorepaso[0.62\linewidth]{respuesta-joroba}{Subida inesperada: $Y$ en joroba; precios después}

% ==========================================================================
\bloque[9]{II. Rigideces y fricciones financieras\label{rep:II}}
\begin{ideaclave}
Basta que los precios no se ajusten todos a la vez para que mover el tipo nominal mueva el real y la demanda. Los contratos salariales alargan el efecto; las fricciones financieras lo amplifican; los hogares distintos cambian el canal principal: de la sustitución intertemporal a la renta y el empleo.
\end{ideaclave}

\sub{II.1. Rigideces de precios: el modelo neokeynesiano\label{rep:II.1}}
\begin{itemize}
\item \emph{Idea}: con precios flexibles $i$↑ va a precios y $r$ no cambia; con revisión gradual, $i$↑ = $r$↑. Por qué rígidos \vertema{3.A.7}: \concepto{costes de menú}: no cambiar el precio cuesta poco a la empresa, pero la suma tiene efectos grandes (\autor{Mankiw}, 1985; \autor{Akerlof}-\autor{Yellen}, 1985); \concepto{rigideces reales} lo amplían (\autor{Ball}-\autor{Romer}, 1990). \emph{Pyme}: tarifa 1-2 veces al año; ajusta pedidos.
\item \emph{Supuestos}: hogar representativo, aversión $\theta$ (EIS $1/\theta$), $\beta$; competencia monopolística, \autor{Calvo} (1983): mantiene el precio con prob.\ $\eta$, dura $1/(1-\eta)$; sin capital, sector público ni exterior; $x_t$ = brecha frente a precios flexibles. 3 ecuaciones, que se cierran entre sí (regla $\to i$, IS $\to x$, NKPC $\to\pi$):
\begin{formula}
\centering $\displaystyle\begin{aligned}
x_t&=\mathbb E_t[x_{t+1}]-\tfrac1\theta\big(i_t-\mathbb E_t[\pi_{t+1}]-r^n_t\big)\\
\pi_t&=\beta\,\mathbb E_t[\pi_{t+1}]+\kappa x_t+u_t\\
i_t&=\rho+\phi_\pi\pi_t+\phi_x x_t+\varepsilon_t
\end{aligned}$
\end{formula}
\item[] $\kappa$↓ si $\eta$↑; $u_t$ costes (\autor{Clarida}-\autor{Galí}-\autor{Gertler}, 1999); $\phi_\pi>1$ (\emph{principio de \autor{Taylor}}); \autor{Taylor} (1993): $i=1+1{,}5\pi+0{,}5x$.
\item \emph{Clave}: IS hacia delante, la brecha depende de toda la senda esperada del tipo real (lectura a la \autor{Wicksell}: neowickselliano, \autor{Woodford}, 2003):
\funcionclave{x_t=-\frac1\theta\sum_{k\geq0}\mathbb E_t\big[i_{t+k}-\pi_{t+k+1}-r^n_{t+k}\big]}
\item \emph{Subida inesperada} $\varepsilon_t>0$: $r$↑, $x$↓, $\pi$↓; en $(x,\pi)$, la demanda a la izquierda, baja por la \autor{Phillips}. \emph{Ojo}: el plano $(x,\pi)$ es aproximado: exige linealizar, perturbaciones AR(1), solución $x_t=\psi_x\varepsilon_t$ y expectativas proporcionales ($\mathbb E_tx_{t+1}=\rho_\varepsilon x_t$). Con $\varepsilon_t=\rho_\varepsilon\varepsilon_{t-1}+\dots$ (\autor{Galí}, 2015):
\begin{formula}
\centering $\displaystyle x_t=-(1-\beta\rho_\varepsilon)\Lambda\varepsilon_t,\quad \pi_t=-\kappa\Lambda\varepsilon_t$
\end{formula}
\item[] Precios de 1 año, $\rho_\varepsilon=0{,}5$, $\varepsilon$ de 1 punto anual: $r$ +0,52; $x$ $-0{,}26$\,\%; $\pi$ $-0{,}35$; $i$ solo +0,34 (la regla responde a la caída).
\item \emph{Implicaciones}: no neutral a CP, no cambia la $Y$ natural; más rígidos ($\kappa$↓) $\Rightarrow$ más efecto en $Y$, menos en $\pi$; canal = tipo real (sustitución intertemporal); \emph{expectativas}: más persistente, más contractiva, aun con $i$ \emph{más bajo} ($\rho_\varepsilon=0{,}8$: $i$ $-0{,}51$, $r$ +0,22): $i$ mal indicador del tono. \concepto{Divina coincidencia}: demanda, sin dilema; costes $u_t$, dilema (\autor{Goodfriend}-\autor{King}, \autor{Rotemberg}-\autor{Woodford}, 1997; nombre: \autor{Blanchard}-\autor{Galí}, 2007).
\item \emph{Ext.}: reglas por estado: \emph{efecto selección}, casi neutral (\autor{Golosov}-\autor{Lucas}, 2007); multiproducto, vuelve a \autor{Calvo} (\autor{Midrigan}, 2011). \autor{Phillips} híbrida (\autor{Galí}-\autor{Gertler}, 1999); información rígida (\autor{Mankiw}-\autor{Reis}, 2002). Medianos: CEE (2005), \autor{Smets}-\autor{Wouters} (2007).
\item \emph{Evid.}: EE.\,UU., mitad de precios $<$4,3 meses con rebajas (\autor{Bils}-\autor{Klenow}, 2004), más sin ellas (\autor{Nakamura}-\autor{Steinsson}, 2008); zona del euro $\approx$8\,\% al mes sin rebajas, sube con la inflación de 2022 (depende del estado; \autor{Gautier} et al.). Para la joroba, salarios rígidos.
\item \emph{Val.}: optimización + rigideces, bienestar. \emph{Contra}: todo por la EIS (baja); sin joroba, bancos ni hogares distintos.
\end{itemize}
\pareja{nek-perturbacion-monetaria}{Subida inesperada en el NEK: DA a la izquierda}{fischer-taylor}{Contratos, $a=0{,}1$: \autor{Fischer} 2 periodos; \autor{Taylor} decae a $\rho_w\approx0{,}52$}

\sub{II.2. Rigideces de salarios: contratos escalonados\label{rep:II.2}}
\begin{itemize}
\item \emph{Idea}: respuesta a la inefectividad: salarios fijados antes de las perturbaciones y no todos a la vez $\Rightarrow$ efectos reales con expectativas racionales. \emph{Pyme}: convenio de 3 años previo a 2022.
\item \emph{Supuestos}: logaritmos; $y_t=m_t-p_t$; contratos de 2 periodos, la mitad cada periodo (\emph{escalonados}); $p$ = media de salarios; deseado $p^*_t=p_t+a\,y_t$, $0<a\leq1$ ($a$ pequeña = rigidez real).
\item \textbf{\autor{Fischer}} (1977), \emph{predeterminados}: $w^j_t=\mathbb E_{t-j}[p^*_t]$, $p_t=\frac12(w^1_t+w^2_t)$:
\begin{formula}
\centering $\displaystyle y_t=\big(m_t-\mathbb E_{t-1}m_t\big)+\tfrac{1}{1+a}\big(\mathbb E_{t-1}m_t-\mathbb E_{t-2}m_t\big)$
\end{formula}
\item[] 1.er término: sorpresa (\autor{Lucas}); 2.º, nuevo: lo sabido en $t-1$, tras firmar en $t-2$, mueve $Y$.
\item \textbf{\autor{Taylor}} (1980), \emph{fijos} en el contrato: $w^c_t=\frac12(p^*_t+\mathbb E_t p^*_{t+1})$, $p_t=\frac12(w^c_t+w^c_{t-1})$; $m$ paseo aleatorio:
\begin{formula}
\centering $\displaystyle w^c_t=\rho_w w^c_{t-1}+(1-\rho_w)m_t,\quad \rho_w=\frac{1-\sqrt a}{1+\sqrt a}$
\end{formula}
\item[] Subida permanente de $m$: $y_k=\frac12(1+\rho_w)\rho_w^{\,k}$. \emph{Forma}: \autor{Fischer}, escalón (1 sorpresa; $1/(1+a)$: la mitad que firma sube poco porque la otra mantiene $p$; 0 cuando todos firmaron después); \autor{Taylor}, caída geométrica a la tasa $\rho_w$ desde $\frac12(1+\rho_w)$. \autor{Calvo} ($\eta=0{,}5$, $\beta=1$, $a=0{,}1$): $y_k=0{,}8^{k+1}$, más persistente.
\item \emph{Implicaciones}: \emph{política sistemática estabiliza}: actúa \emph{después} de firmados los contratos, sin saber más que el público; \emph{persistencia}: \autor{Fischer} dura lo que el contrato (1; 0,91; 0); \autor{Taylor}, cada grupo mira al otro: $a=0{,}1$, $\rho_w\approx0{,}52$: 0,76; 0,40; 0,21; 0,11…; dura más que los contratos; inercia: la rigidez real amplía la nominal.
\item \emph{Ext.}: \autor{Calvo} (1983), probabilidad en vez de duración: sin memoria, sale de optimizar, se agrega en la NKPC, da el coste de la inflación (dispersión); \autor{Rotemberg} (1982), coste cuadrático, misma curva. Salarios a la \autor{Calvo}: 2 \autor{Phillips}, estabilizar solo precios no es óptimo (\autor{Erceg}-\autor{Henderson}-\autor{Levin}, 2000). Salario real rígido: rompe la divina coincidencia (\autor{Blanchard}-\autor{Galí}, 2007; con paro, 2010) \vertema{3.A.28}. Indexación óptima parcial (\autor{Gray}, 1976).
\item \emph{Evid.}: Europa: salarios $\approx$15 meses sin cambio (precios $\approx$10), en enero (\autor{Druant} et al., 2009); pico en cero, pocas bajadas (\autor{Dickens} et al., 2007): paro de la periferia tras 2008 (\autor{Schmitt-Grohé}-\autor{Uribe}, 2016), razón contra $\pi=0$ (\autor{Akerlof}-\autor{Dickens}-\autor{Perry}, 1996) \vertema{3.A.41}; fricción más importante en los medianos.
\item \emph{Val.}: puente expectativas racionales-rigideces, persistencia. \emph{Contra}: contratos exógenos; ¿por qué no indexar? (\autor{Barro}, 1977).
\end{itemize}

\sub{II.3. Fricciones financieras: los canales del crédito\label{rep:II.3}}
\begin{itemize}
\item Sin fricciones, \autor{Modigliani}-\autor{Miller}: la financiación no importa \vertema{3.A.34}; con información asimétrica, sí \vertema{3.A.13}. \emph{Fundamento}: racionamiento (\autor{Stiglitz}-\autor{Weiss}, 1981); verificación costosa (\autor{Townsend}, 1979) $\Rightarrow$ \concepto{prima de financiación externa}, ↓ con el patrimonio (\autor{Bernanke}-\autor{Gertler}, 1989). \autor{Bernanke} (1983): Gran Depresión por quiebras bancarias.
\item \concepto{Canal del crédito bancario} (\autor{Bernanke}-\autor{Blinder}, 1988, 1990; antes \autor{Blinder}-\autor{Stiglitz}, 1983): préstamos y bonos no sustitutos perfectos $\Rightarrow$ la política mueve la \emph{oferta} de crédito; por la oferta de préstamos, no por el multiplicador. \emph{Evid.}: bancos pequeños y menos líquidos recortan más (\autor{Kashyap}-\autor{Stein}, 2000); España: bancos con menos capital o liquidez conceden menos y la empresa rechazada no lo encuentra en otro (\autor{Jiménez} et al., 2012).
\item \concepto{Acelerador financiero} (\autor{Bernanke}-\autor{Gertler}-\autor{Gilchrist}, 1999): NEK con capital; empresarios compran capital ($q_t$) con riqueza neta RN y deuda:
\funcionclave{\mathbb E_t[R^K_{t+1}]=s\Big(\frac{\mathrm{RN}_{t+1}}{q_tK_{t+1}}\Big)R_{t+1},\quad s'<0}
\item[] \emph{Doble golpe}: sube el tramo plano ($r\to r'$) y cae RN (intereses variables, flujo de caja, garantías $q_t$): $I_0\to I_1$, más que solo con $r$. Bajada: $\downarrow i\Rightarrow\uparrow$flujo, garantías, $P$ no anticipado $\Rightarrow\uparrow$RN$\Rightarrow\downarrow s\Rightarrow\uparrow I,C$.
\item \emph{Implicaciones}: amplificación y persistencia (RN tarda en reconstruirse); \emph{asimetría} (recesiones, pequeñas); deflación por deudas (\autor{Fisher}, 1933). \emph{Pyme}: menos pedidos, local que vale menos: el banco pide más (el banco presta el paraguas cuando hace sol y lo quita cuando llueve).
\item \emph{Ext.}: garantías (\autor{Kiyotaki}-\autor{Moore}, 1997) \vertema{3.A.42}; capital de los bancos (\autor{Gertler}-\autor{Kiyotaki}, 2010; \autor{Gertler}-\autor{Karadi}, 2011): puente a las compras.
\item \emph{Evid.}: ventas y existencias de pequeñas caen más (\autor{Gertler}-\autor{Gilchrist}, 1994); subida sorpresa sube diferenciales (\autor{Gertler}-\autor{Karadi}, 2015). \emph{Val.}: amplificación, potencia en recesiones con deuda; \emph{contra}: lineales, explican mejor la caída que el auge.
\end{itemize}
\graficorepaso[0.62\linewidth]{coste-fondos-politica}{Doble golpe: $r\to r'$ y RN$\to$RN$'$: $I_0\to I_1$}

\sub{II.4. Hogares distintos: el canal redistributivo\label{rep:II.4}}
\begin{itemize}
\item \emph{Idea}: EIS baja; muchos viven al día (\emph{hand-to-mouth}): EE.\,UU.\ $\approx$1/3, la mayoría con riqueza ilíquida (\autor{Kaplan}-\autor{Violante}-\autor{Weidner}, 2014): el tipo les llega por renta y empleo. \emph{Familia}: cuota y empleo.
\item \concepto{HANK} (\autor{Kaplan}-\autor{Moll}-\autor{Violante}, 2018): hogares distintos en renta, riqueza y liquidez, riesgos no asegurables, límite al endeudamiento, activo líquido e ilíquido \vertema{3.A.29}:
\begin{formula}
\centering $\displaystyle \mathrm dC_0=\underbrace{\sum_t\tfrac{\partial C_0}{\partial r_t}\mathrm dr_t}_{\text{directo}}+\underbrace{\sum_t\Big(\tfrac{\partial C_0}{\partial(wL)_t}\mathrm d(wL)_t+\tfrac{\partial C_0}{\partial T_t}\mathrm dT_t\Big)}_{\text{indirectos}}$
\end{formula}
\item[] RANK: casi todo directo; HANK: indirectos $\approx$4 de cada 5 €.
\item \concepto{Canal redistributivo} (\autor{Auclert}, 2019): reparto entre propensiones distintas: \emph{renta} (paro de rentas bajas), \emph{\autor{Fisher}} ($\pi$ inesperada↓: de deudores a acreedores), \emph{exposición a tipos} (ganan depósitos, pierden deudas variables). TANK (\autor{Galí}-\autor{López-Salido}-\autor{Vallés}, 2007; \autor{Bilbiie}, 2008, 2020): amplifica si la renta de los que viven al día es más cíclica.
\item \emph{Impl.}: actúa por renta y empleo; España (variable) más fuerte; quién paga intereses de la deuda pública importa \vertema{3.A.40}; promesas lejanas pesan menos (\autor{McKay}-\autor{Nakamura}-\autor{Steinsson}, 2016); redistribuye, pero redistribuir es de la fiscal.
\item \emph{Evid.}: Noruega, indirectos dominan con el tiempo, endeudados recortan más (\autor{Holm}-\autor{Paul}-\autor{Tischbirek}, 2021); RU y EE.\,UU.: con hipoteca e inquilinos ajustan, propietarios sin deuda no (\autor{Cloyne}-\autor{Ferreira}-\autor{Surico}, 2020). \emph{Val.}: encajan con microdatos; dependen de la calibración.
\end{itemize}

% ==========================================================================
\bloque[9]{III. Transmisión de la política no convencional\label{rep:III}}
\begin{ideaclave}
En el límite, el NEK da una espiral deflacionista y sus salidas: bajar el suelo, subir la inflación esperada y bajar las primas. Son los canales no convencionales; funcionaron, con menos precisión que el tipo oficial, y deshacerlos movió menos los mercados que hacerlos.
\end{ideaclave}

\sub{III.1. El límite inferior de los tipos\label{rep:III.1}}
\begin{itemize}
\item Efectivo rinde 0\,\% $\Rightarrow$ \concepto{límite inferior efectivo} $i^{\mathrm{ELB}}$ (ZLB; ELB si algo negativo) \vertema{3.A.35}; suelo del tipo real $r_t\geq i^{\mathrm{ELB}}-\mathbb E_t\pi_{t+1}$. Fed 0-0,25\,\% 2008-2015 y 2020-2022; BCE negativo jun.\ 2014-jul.\ 2022 (mín.\ $-0{,}50$\,\%).
\item \emph{Por qué}: $r^n$ cae 30 años: \concepto{estancamiento secular} (\autor{Hansen}, 1939; \autor{Summers}, 2013-2014): ahorro↑ (envejecimiento, precaución, desigualdad), inversión↓ (población, productividad, capital barato); $r^n$ $-300$ pb (\autor{Rachel}-\autor{Summers}, 2019). \autor{Blanchard}: 2 umbrales, $r<\gamma_Y$ y el ELB; papel mayor de la fiscal. 2023: \autor{Summers}, terminó; \autor{Blanchard}, no \vertema{3.A.40}.
\end{itemize}
\graficorepaso[0.6\linewidth]{estancamiento-secular}{Estancamiento secular: $r^n$ bajo el límite, sobra ahorro}

\sub{III.2. El modelo neokeynesiano en el límite\label{rep:III.2}}
\begin{itemize}
\item \emph{Idea}: $r^n$ bajo el suelo $\Rightarrow$ $x$↓, $\pi$↓, $\pi^e$↓, $r$↑: espiral deflacionista (\autor{Krugman}, 1998; \autor{Eggertsson}-\autor{Woodford}, 2003) \vertema{3.A.41}; no es la trampa de \autor{Keynes}. \emph{Familia}: Euríbor $<0$ 2016-2021.
\item \emph{Supuestos}: 3 ecuaciones; $i_t=\max\{i^{\mathrm{ELB}},\rho+\phi_\pi\pi_t+\phi_xx_t\}$; prima $\omega_t$ en la IS ($i_t+\omega_t$); $r^n_L$ persiste con prob.\ $p$; $\mathbb E_tx_{t+1}=px_L$, $\mathbb E_t\pi_{t+1}=p\pi_L$.
\item \emph{Desarrollo}: oferta $\pi_L=\frac{\kappa}{1-\beta p}x_L$; demanda en el límite $\pi_L=\frac{\theta(1-p)}{p}x_L-\frac{r^n_L-\omega_L-i^{\mathrm{ELB}}}{p}$, creciente (codo); equilibrio si $\theta(1-p)(1-\beta p)>p\kappa$:
\funcionclave{x_L=\frac{(1-\beta p)(r^n_L-\omega_L-i^{\mathrm{ELB}})}{\theta(1-p)(1-\beta p)-p\kappa}<0}
\item[] $\theta=1$, $\beta=0{,}99$, $\kappa=0{,}05$, $p=0{,}7$, $r^n$ 1 punto bajo el suelo: $x$ $-1{,}3$\,\%, $\pi$ $-0{,}9$; con $\kappa=0{,}1$, $-3{,}5$\,\%: \emph{paradoja de la flexibilidad} (\autor{Eggertsson}-\autor{Krugman}, 2012).
\item \emph{Implicaciones (canales)}: \textbf{tipos negativos} ($i^{\mathrm{ELB}}$↓; tipo de reversión, \autor{Abadi}-\autor{Brunnermeier}-\autor{Koby}, 2023). \textbf{Orientación} ($\mathbb E\pi$↑: promesa de tipos bajos tras la perturbación; BCE 4/7/2013; «prometer de forma creíble que será irresponsable», \autor{Krugman}; délfica vs.\ odiseica; \concepto{paradoja de la orientación}, \autor{McKay}-\autor{Nakamura}-\autor{Steinsson}, 2016). \textbf{Compras} ($\omega$↓; sin primas irrelevantes, \autor{Wallace}, 1981; señalización, carteras/hábitat, liquidez, crédito, cambio y riqueza; APP 2014, PEPP 2020). \textbf{Fiscal}: multiplicador mayor \vertema{3.A.40}.
\item \emph{Ext.}: préstamos condicionados (TLTRO 2014; +186 M€ por 1.000 M€, \autor{Andrade} et al.); protección de la transmisión (OMT 2012, TPI 2022); control de la curva (Japón 2016-2024); dinero helicóptero. \emph{Comodín}: suprimir el efectivo (\autor{Rogoff}, 2016; España, 57\,\% de compras en tienda en efectivo, 2024) y CBDC remunerada; euro digital sin intereses, piloto 2027, emisión 2029.
\item \emph{Evid.}: orientación, tipos a 1-2 años; compras, los largos (\autor{Swanson}, 2021), en parte por la señal (\autor{Krishnamurthy}-\autor{Vissing-Jorgensen}, 2011); negativos: menos crédito en bancos con depósitos (\autor{Heider} et al., 2019); tipo sombra zona del euro $-7{,}8$\,\% (ago.\ 2019).
\item \emph{Val.}: ordena las salidas; \emph{contra}: exagera las promesas lejanas y sin fricciones no da papel a las compras. \emph{Salida} asimétrica: reducir el balance subió los largos mucho menos (\autor{Du}-\autor{Forbes}-\autor{Luzzetti}, 2024).
\end{itemize}
\pareja{nek-limite-inferior}{DA con codo: $E_0$ en la espiral; salidas $\to E_1$}{no-convencionales-canales}{Salidas $\to$ términos del modelo $\to$ canales}

% ==========================================================================
\bloque[1]{Conclusión}
\begin{itemize}
\item \textbf{Recapitulación}: «En conclusión,» varios canales a la vez (tipo, activos, cambio, crédito, expectativas) porque hay fricciones: precios y salarios lentos, patrimonio que importa; crean y amplifican la transmisión; con hogares distintos cambia a quién llega; en el límite, palabra y balance.
\item \textbf{Relevancia}: España (variable, bancos): más rápida y desigual (2022-2023).
\item \textbf{Extensiones}: \vertema{3.A.36}, \vertema{3.A.35}, \vertema{3.A.7}, \vertema{3.A.40}, \vertema{3.A.39}, \vertema{3.A.41}, \vertema{3.A.42}, \vertema{3.B.14}, \vertema{3.B.19}, \vertema{3.B.43}; euro digital remunerado: canal directo al público.
\item \textbf{Opinión}: consenso en la joroba con retraso; menos acuerdo en los canales de las compras: bajaron los largos, pero tuvieron que ser muy grandes (\autor{Walsh}, 2017).
\item \textbf{Idea final}: las cosas de palacio van despacio: actúa hoy sobre la economía de dentro de 2 años; mirar adelante y ser creíble, porque cuando el tipo ya no baja, su palabra es lo que le queda.
\end{itemize}

% ==========================================================================
\bloque{Autores y fechas}
\hito{1752}{\autor{Hume}: identidad del gasto; 1911, \autor{Fisher}: teoría cuantitativa}
\hito{1936}{\autor{Keynes}: canal del tipo de interés}
\hito{1972}{\autor{Lucas}: islas; 1975, \autor{Sargent}-\autor{Wallace}: inefectividad}
\hito{1977}{\autor{Fischer}: contratos predeterminados}
\hito{1980}{\autor{Taylor}: contratos escalonados}
\hito{1981}{\autor{Stiglitz}-\autor{Weiss}: racionamiento; \autor{Wallace}: irrelevancia del balance}
\hito{1983}{\autor{Calvo}; \autor{Bernanke}: Gran Depresión}
\hito{1988}{\autor{Bernanke}-\autor{Blinder}: crédito bancario; 1989, \autor{Bernanke}-\autor{Gertler}: prima}
\hito{1995}{\autor{Bernanke}-\autor{Gertler}: canales del crédito}
\hito{1996}{\autor{Christiano}-\autor{Eichenbaum}-\autor{Evans}: VAR, joroba (\autor{Sims}, 1992: rompecabezas)}
\hito{1998}{\autor{Krugman}: «irresponsable» creíble}
\hito{1999}{\autor{Bernanke}-\autor{Gertler}-\autor{Gilchrist}: acelerador}
\hito{2003}{\autor{Eggertsson}-\autor{Woodford}: orientación en el límite}
\hito{2007}{\autor{Blanchard}-\autor{Galí}: divina coincidencia; \autor{Smets}-\autor{Wouters}}
\hito{2015}{\autor{Gertler}-\autor{Karadi}: alta frecuencia}
\hito{2016}{\autor{McKay}-\autor{Nakamura}-\autor{Steinsson}; \autor{Rogoff}}
\hito{2018}{\autor{Kaplan}-\autor{Moll}-\autor{Violante}: HANK}
\hito{2019}{\autor{Auclert}: canal redistributivo}
\hito{2023}{\autor{Abadi}-\autor{Brunnermeier}-\autor{Koby}: tipo de reversión}

% ==========================================================================
\bloque{Datos clave}
\begin{itemize}
\item Euríbor 12 m: \dato{$-0{,}50\to4{,}16$\,\%}{BCE, 2021-2023}; inflación: \dato{10,6\,\%}{Eurostat, oct.\ 2022}; \dato{4\,\%}{oct.\ 2026}; depósito \dato{2,50\,\%}{BCE, sept.\ 2026}.
\item Efecto máximo en precios: \dato{29 meses}{Havranek-Rusnak, 2013}; precios bajan tras \dato{22 meses}{Romer-Romer, 2004}; $Y$ industrial \dato{$-4{,}3$\,\%}{Romer-Romer, 2004}.
\item Traspaso 2022-23: empresas \dato{+350 pb}{BCE MIR}; hipotecas \dato{$>$200 pb}{BCE MIR}; vista \dato{0,4 / 0,2\,\%}{zona / España}.
\item Hipotecas vivas: España \dato{1,2 $\to$ 3,7\,\%}{BCE MIR, jun.\ 2022-feb.\ 2024}; zona \dato{1,6 $\to$ 2,5\,\%}{BCE MIR}.
\item Precios que cambian al mes (sin rebajas): \dato{$\approx$8\,\%}{zona del euro, Gautier et al.}; salarios fijos \dato{$\approx$15 meses}{Druant et al., 2009}.
\item Hogares que viven al día, EE.\,UU.: \dato{$\approx$1/3}{Kaplan-Violante-Weidner, 2014}; efectos indirectos \dato{$\approx$4/5}{Kaplan-Moll-Violante, 2018}.
\end{itemize}

% ==========================================================================
\bloque{Preguntas del tribunal}
\begin{enumerate}
\item \textbf{¿Por qué tiene efectos reales la política monetaria?} Rigideces nominales: mover $i$ mueve $r$; reales y fricciones financieras amplían; a LP, neutral.
\item \textbf{¿Qué refuta \autor{Fischer} de la inefectividad?} Con contratos firmados antes, la regla actúa tras la firma y mueve $Y$ aun con expectativas racionales: la clave era la flexibilidad.
\item \textbf{¿Por qué un tipo nominal más bajo puede ser contractivo?} Con perturbación persistente la inflación cae tanto que $r$ sube ($\rho_\varepsilon=0{,}8$: $i$ $-0{,}51$, $r$ +0,22).
\item \textbf{¿Qué añade el acelerador financiero?} Doble golpe: $r$↑ y RN↓ $\Rightarrow$ prima↑; amplificación, persistencia, asimetría.
\item \textbf{¿Por qué el HANK cambia el canal?} EIS baja y hogares sin liquidez: $\approx$4/5 del efecto es indirecto (renta, empleo).
\item \textbf{¿Por qué funcionan las compras si \autor{Wallace} dice que no?} Mercados imperfectos ($\omega>0$): señalización, carteras/hábitat, liquidez, crédito, cambio.
\item \textbf{¿Qué salidas da el modelo en el límite?} Bajar $i^{\mathrm{ELB}}$, subir $\mathbb E\pi$ (orientación), bajar $\omega$ (compras); fuera del modelo, crédito, primas soberanas, curva; comodín, quitar el efectivo.
\end{enumerate}

\begin{trampa}
\begin{itemize}
\item \autor{Stanley Fischer} (contratos, 1977) $\neq$ \autor{Irving Fisher} (ecuación de cambio, 1911; deflación por deudas, 1933). Regla del $k$\,\% $\neq$ regla de \autor{Friedman} ($i=0$).
\item Canal del crédito = préstamo bancario (\autor{Bernanke}-\autor{Blinder}, 1988) + balance (\autor{Bernanke}-\autor{Gertler}, 1995); \autor{Blinder}-\autor{Stiglitz} (1983), antecedente.
\item Divina coincidencia: resultado de \autor{Goodfriend}-\autor{King} y \autor{Rotemberg}-\autor{Woodford} (1997); nombre de \autor{Blanchard}-\autor{Galí} (2007); perturbación de costes, \autor{Clarida}-\autor{Galí}-\autor{Gertler} (1999).
\item Identidad $MV\equiv PT$ (\autor{Hume}) $\neq$ teoría cuantitativa (\autor{Fisher}: identidad + supuestos causales).
\item \autor{Calvo}: duración $1/(1-\eta)$, $\eta$ = prob.\ de \emph{no} revisar; \autor{Taylor}: más rigidez real, más persistencia.
\item Paradoja de la orientación: en el modelo la promesa lejana tiene tanto efecto como la cercana.
\item Límite inferior no es exactamente 0 (BCE $-0{,}50$\,\%); bajo el tipo de reversión, contrae el crédito.
\item \autor{Gertler}-\autor{Karadi} (2015), EE.\,UU.; BCE: \autor{Altavilla} et al., \autor{Jarociński}-\autor{Karadi}.
\end{itemize}
\end{trampa}

\end{columnas}
\end{document}
